\scanpage{87}
par produit tensoriel \unsure{des|de} formes [N.B. le produit tensoriel de deux formes \deleted{symétriques, ou de} de $=$ type de symétrie est symétrique \textendash\ de deux types opposés de symétrie est antisymétrique]. Bien entendu, ils ne sont pas en général bijectifs (mais du moins injectifs).

On a aussi
\[
\unsure{\Phi|\mathscr{L}|\mathcal{L}}^s(M'\times M'')\simeq
\Phi^s(M')\times\Phi^s(M'')\times B(M',M'')
\]
où
\[
B(M',M'')=\operatorname{Hom}(M'\otimes M'',\mathbb{Q}(p)).
\]

Ceci posé, on \deleted{veut pour \unsure{dt|tt}} \inserted{a dans} tt $\Phi^s(M)$ un cône \inserted{convexe} \deleted{\unsure{f}}fermé $\Phi^+(M)$.

On a les propriétés suivantes :

(i) $\Phi^+(M)$ a un intérieur non vide.

(ii) Un élément de $\Phi^+(M)$ est dans l'int. de $\Phi^+(M)$ ssi il est non dég.

\margin{formes $\in\Phi^+(M)$ sont dites pseudo-amples, celles dans l'intérieur sont dites \uline{amples}}
% slanted

(iii) Pour $M\to N$, $\Phi^+(N)$ s'envoie dans $\Phi^{\unsure{+|e}}(M)$.

(iv) \deleted{\unsure{(iii)} \unsure{Pour}} Si $M\longrightarrow N$ injectif, \deleted{\(\Phi^+(N)\)} l'intérieur de $\Phi^+(N)$ s'envoie dans l'intérieur de $\Phi^+(M)$ (formes amples $\longrightarrow$ formes amples)
% archive sheet number: 86
